\chapter{補遺}


\setcounter{chapter}{0}
\newtheorem{thm2}{定理}[chapter]

採用した公理系および公準の完全性について幾つか述べる。
ここでいう公理系の完全性とは、公理系の定める言語における任意の閉論理式（自由変数を含まない論理式）について、必ず真偽が定まることである。

\ref{261008235001} の公理系（公準を含まない）は、実数体の公理系において上限または下限の存在の公理かデデキント切断の公理になっているはずの部分が、集合を用いずに一階述語論理で表したデデキント切断の公理になっている。
ここでは、この公理を「一階デデキント切断の公理」と呼ぶことにしよう。
この公理系は、順序体の公理系に一階デデキント切断の公理を加えたものとも言える。

% 定理 I

\begin{thm2}
  一階デデキント切断の公理を持つ順序体の公理系のモデルはどれも、一階述語論理の任意の閉論理式に対して全く同じ真偽を割り当てる。
\end{thm2}

\begin{proof}
  そのような公理系においては、一階述語論理で表されるあらゆるデデキント切断が存在する。

  実数体におけるデデキント切断は、条件を満たす二つの集合が
  \[\{\ldots s_{l+1},\ldots s_{l+n},\ldots,r\},\ \{\ldots t_{m+1},\ldots t_{m+n},\ldots\}\]
  か
  \[\{\ldots s_{l+1},\ldots s_{l+n},\ldots\},\ \{r,\ldots t_{m+1},\ldots t_{m+n},\ldots\}\]
  の形であるが、問題になっている公理系においては、
  \[\{x\mid\varphi(x)\},\ \{y\mid\psi(y)\}\]
  という、論理式 \(\varphi(x),\ \psi(y)\) を用いて表されるものに限られる。

  したがって、ここではモデルが一階述語論理の論理式に与える真偽についてのみ考えているのだから、一階述語論理で表せないものを無視すれば、言い表すことのできるデデキント切断の境界となる点を最大限含むという意味で、当の公理系のモデルは一つしかない。
  もちろん、集合で考えれば話は全く別で、無数にあることになる。

  しかし、ここまでの大前提は実は「（集合を用いた）実数体の公理系のモデルが本質的に一つだったら」ということであった。
  ゆえに、\kenten{実数体の公理系の範疇性}(categoricity)\kenten{により}、定理の結論が言える。
\end{proof}

\begin{col*}
  一階デデキント切断の公理を持つ順序体の公理系は完全である。
\end{col*}

% 定理 II

\begin{thm2}
  一階デデキント切断の公理を持つ順序体の公理系に座標系の公準を付け加えた公理系は、完全である。
\end{thm2}

\begin{proof}
  任意の順序体を対象領域として、そこに \(n\) 次元ベクトルという単に順序付けられた要素の組をわざと違ったやり方で導入して、一見異なる \(n\) 次元座標を二つ構築しても、等号・和・定数倍・内積も含めて、それら二つの間の同型性(isomorphism)は容易に示すことができる（原子論理式を経由する同型の矢印を考えよ）。

  つまり、それらモデルの拡張による新たな論理式への真理値の割り当て方は、一種類しかない。

  ゆえに、一階述語論理の任意の閉論理式に与える真理値のみで判定することにすれば、一階デデキント切断の公理を持つ順序体の公理系のモデルは相互に区別がつかないから、それぞれに座標系を付け加えたものもまた相互に区別がつかない。
  なぜなら、もし異なるとすると、一階述語論理の範囲に縮小した一つのモデルに対して、同型でない二つの座標系が構築可能ということになるからである。

  ゆえに、それらモデルの同値性（どのモデルも閉論理式に与える真理値に関して同等であるということ）は揺るがない。
\end{proof}
